\pdfoutput=1
% ================================================================
% SCX — Machine Translation as Cross-Language Gauge Fixing
% ================================================================
% Compile: pdflatex -interaction=nonstopmode main.tex
% ================================================================
\documentclass[12pt,a4paper]{article}

% ================================================================
% Language & Encoding
% ================================================================
\usepackage[english]{babel}
\usepackage[T1]{fontenc}

% ================================================================
% Page Geometry
% ================================================================
\usepackage[a4paper,top=2cm,bottom=2cm,left=2.5cm,right=2.5cm]{geometry}

% ================================================================
% Math Core
% ================================================================
\usepackage{amsmath,amssymb,amsthm}
\usepackage{mathtools}
\usepackage{bm}
\usepackage{bbm}

% ================================================================
% Graphics & Color
% ================================================================
\usepackage{graphicx}
\usepackage{tikz}
\usepackage{xcolor}
\usepackage{float}

% ================================================================
% Tables & Lists
% ================================================================
\usepackage{booktabs}
\usepackage{enumitem}
\usepackage{paralist}

% ================================================================
% Algorithms
% ================================================================
\usepackage{algorithm}
\usepackage{algpseudocode}

% ================================================================
% Hyperlinks & Cross-References
% ================================================================
\usepackage[colorlinks=true, citecolor=blue, urlcolor=blue]{hyperref}
\usepackage{cleveref}

% ================================================================
% Theorem Environments (English)
% ================================================================
\newtheorem{theorem}{Theorem}[section]
\newtheorem{lemma}[theorem]{Lemma}
\newtheorem{proposition}[theorem]{Proposition}
\newtheorem{corollary}[theorem]{Corollary}
\newtheorem{definition}[theorem]{Definition}
\newtheorem{conjecture}[theorem]{Conjecture}
\newtheorem{assumption}[theorem]{Assumption}

\theoremstyle{remark}
\newtheorem{remark}[theorem]{Remark}
\newtheorem{example}[theorem]{Example}
\newtheorem{protocol}[theorem]{Protocol}

% ================================================================
% Math Shortcuts — Blackboard Bold
% ================================================================
\newcommand{\R}{\mathbb{R}}
\newcommand{\C}{\mathbb{C}}
\newcommand{\N}{\mathbb{N}}
\newcommand{\Z}{\mathbb{Z}}
\newcommand{\Q}{\mathbb{Q}}
\newcommand{\E}{\mathbb{E}}
\newcommand{\Pbb}{\mathbb{P}}

% ================================================================
% Math Shortcuts — Calligraphic
% ================================================================
\newcommand{\D}{\mathcal{D}}
\newcommand{\F}{\mathcal{F}}
\newcommand{\G}{\mathcal{G}}
\newcommand{\T}{\mathcal{T}}
\newcommand{\loss}{\mathcal{L}}
\newcommand{\M}{\mathcal{M}}
\newcommand{\Hilb}{\mathcal{H}}
\newcommand{\A}{\mathcal{A}}
\newcommand{\B}{\mathcal{B}}
\newcommand{\I}{\mathcal{I}}
\newcommand{\J}{\mathcal{J}}
\newcommand{\X}{\mathcal{X}}
\newcommand{\K}{\mathcal{K}}
\newcommand{\Sg}{\mathcal{S}}
\newcommand{\U}{\mathcal{U}}
\newcommand{\V}{\mathcal{V}}
\newcommand{\W}{\mathcal{W}}
\newcommand{\Lang}{\mathcal{L}}
\newcommand{\Embed}{\mathcal{E}}
\newcommand{\Vocab}{\mathcal{V}}
\newcommand{\Trans}{\mathcal{T}}
\newcommand{\Gauge}{\mathcal{G}}

% ================================================================
% SCX Framework Macros
% ================================================================
\newcommand{\SCX}{\textsf{SCX}}
\newcommand{\Yajie}{\textsf{Yajie}}
\newcommand{\Spring}{\textsf{Spring}}
\newcommand{\Cercis}{\textsf{Cercis}}
\newcommand{\Situs}{\textsf{Situs}}
\newcommand{\MoE}{\textsf{MoE}}
\newcommand{\BERT}{\textsf{BERT}}
\newcommand{\XLM}{\textsf{XLM-R}}
\newcommand{\UG}{\textsf{UG}}

% ================================================================
% Operators
% ================================================================
\DeclareMathOperator{\argmax}{arg\,max}
\DeclareMathOperator{\argmin}{arg\,min}
\DeclareMathOperator{\sgn}{sgn}
\DeclareMathOperator{\diag}{diag}
\DeclareMathOperator{\spec}{spec}
\DeclareMathOperator{\softmax}{softmax}
\DeclareMathOperator{\topk}{top\text{-}k}
\DeclareMathOperator{\Cov}{Cov}
\DeclareMathOperator{\Var}{Var}
\DeclareMathOperator{\Tr}{Tr}
\DeclareMathOperator{\BLEU}{BLEU}
\DeclareMathOperator{\CercisScore}{Cercis}
\DeclareMathOperator{\rank}{rank}
\DeclareMathOperator{\supp}{supp}
\DeclareMathOperator{\Dist}{Dist}
\DeclareMathOperator{\TV}{TV}
\DeclareMathOperator{\KL}{KL}
\DeclareMathOperator{\Precision}{Prec}
\DeclareMathOperator{\Recall}{Recall}
\DeclareMathOperator{\BLEUScore}{BLEU}
\DeclareMathOperator{\Align}{Align}
\DeclareMathOperator{\Proj}{Proj}
\DeclareMathOperator{\Translation}{Trans}
\DeclareMathOperator{\LangSpace}{LangSpace}
\DeclareMathOperator{\Invariant}{Inv}
\DeclareMathOperator{\GaugeFix}{GaugeFix}

% ================================================================
% Proof Status Markers
% ================================================================
\newcommand{\rigorFull}{\textbf{[Full Proof]}}
\newcommand{\rigorPartial}{\textbf{[Partial Proof]}}
\newcommand{\rigorSketch}{\textbf{[Proof Sketch]}}
\newcommand{\openproblem}{\textbf{[Open Problem]}}

% ================================================================
% Honest Critique
% ================================================================
\newcommand{\honestcrit}[1]{\textcolor{red}{\textbf{[Honest Critique]} #1}}

% ================================================================
% Document Metadata
% ================================================================
\title{\textbf{Machine Translation as Cross-Language Gauge Fixing:\\
BLEU, Multilingual Models, and Universal Grammar\\
Through the SCX Lens}}
\author{SCX}
\date{\today}

% ================================================================
\begin{document}
\maketitle

\begin{abstract}
We present a unified gauge-theoretic interpretation of machine translation (MT)
within the \SCX{} (State-Conditioned eXpertise) framework. The central thesis is
that each human language constitutes a distinct coordinate system for encoding
semantic content, and translation between languages is precisely a gauge
transformation: meaning is the gauge-invariant content that persists across
all language coordinate charts, while surface form (phonology, morphology,
syntax) is gauge-variant. We establish nine exact correspondences between MT
and \SCX{}. (1) \textbf{Translation as gauge transformation}: source language
$L_s$ to target language $L_t$ preserves semantic content (gauge-invariant)
while changing surface form (gauge-variant). (2) \textbf{BLEU score as Cercis
between reference and hypothesis}: BLEU's n-gram precision is a quality
measure $Q$, and its brevity penalty plays the role of a novelty term $N$;
the BLEU score is structurally identical to the \Cercis{} score $S(s)=Q(s)+
\eta(t)\cdot N(s)$. (3) \textbf{Multilingual models (mBERT, XLM-R) as
multi-expert shared latent space}: the encoder learns a gauge-invariant
representation across languages, with each language-specific layer
performing a gauge fixing. (4) \textbf{Code-switching as gauge ambiguity}:
alternating between languages within a single utterance constitutes undefined
gauge at the switch point, analogous to coordinate singularities. (5)
\textbf{Universal Grammar (UG) as the gauge-invariant substrate}: UG is
the set of linguistic features invariant under all language-to-language
gauge transformations. (6) \textbf{MT evaluation as audit}: human evaluation
(MQM, DQF) constitutes an $M$-expert audit where inter-annotator agreement
is the \Cercis{} score. (7) \textbf{Word embeddings as gauge fields}: the
embedding space $\Embed$ defines a connection between language coordinate
charts, and cross-lingual alignment learns the parallel transport map.
(8) \textbf{Neural MT as learned gauge fixing}: the attention mechanism
learns to optimize the gauge transformation, and the transformer is a
gauge-equivariant architecture. (9) \textbf{Translationese as gauge drift}:
translated texts exhibit systematic deviations from native language due to
incomplete gauge fixing. We derive formal theorems connecting each MT
phenomenon to \SCX{} constructs, prove a BLEU--Cercis isomorphism theorem,
establish the UG invariance principle, and prove that the multilingual
model latent space forms an approximate UG manifold. A complete numerical
verification suite is provided in \texttt{verify\_linguistics.py}.
\end{abstract}


% ================================================================
\section{Introduction}
\label{sec:introduction}
% ================================================================

Machine translation has experienced a revolution in the past decade. From
statistical phrase-based systems through neural sequence-to-sequence models
to large-scale multilingual transformers such as mBERT~\cite{devlin2019bert}
and XLM-R~\cite{conneau2020unsupervised}, the quality of automatic translation
has improved dramatically, approaching human parity for several high-resource
language pairs~\cite{hassan2018achieving}. Yet a unified theoretical framework
that connects translation quality evaluation, multilingual representation
learning, code-switching, and universal grammar remains elusive.

This paper argues that such a framework already exists: the \SCX{} theory of
multi-expert consensus and bias detection. The connection is not merely
metaphorical but \emph{structural}: the same mathematics that describes
how an ensemble of experts converges to a consensus score also describes
how a source sentence in one language is transformed into a target sentence
in another language.

\begin{definition}[Language as Coordinate System]
\label{def:language_coordinate}
A \textbf{language coordinate system} $\Lang_a$ is a tuple
$(\Vocab_a, \Phi_a, \Sigma_a)$ where:
\begin{itemize}[leftmargin=*]
  \item $\Vocab_a$ is the vocabulary (lexical inventory) of language $a$;
  \item $\Phi_a: \Vocab_a^* \to \M$ is the meaning mapping from strings
        of $\Vocab_a$ to an abstract meaning space $\M$;
  \item $\Sigma_a \subset \Vocab_a^*$ is the set of grammatical strings
        (well-formed sentences) of language $a$.
\end{itemize}
A \textbf{semantic content} $m \in \M$ is a point in meaning space. A
sentence $x_a \in \Sigma_a$ is a coordinate representation of $m$ in
the chart $(\Vocab_a, \Phi_a)$.
\end{definition}

The key insight is that the same semantic content can be expressed in
multiple language coordinate systems, just as the same geometric point
can be described by different coordinate charts on a manifold.

\begin{remark}[Coordinate Freedom of Language]
\label{rem:coordinate_freedom}
``Snow'' (English), ``Schnee'' (German), ``nieve'' (Spanish), ``yuki'' (Japanese),
and ``xue'' (Mandarin) all refer to the same semantic content---the frozen
crystalline form of water. The phonemes, graphemes, and morphological
infections are coordinate-dependent; the reference is coordinate-independent.
This is the linguistic manifestation of \emph{gauge freedom}.
\end{remark}

\medskip
\noindent\textbf{Structure of the paper.}
Section~\ref{sec:translation_gauge} formalizes translation as a gauge
transformation. Section~\ref{sec:bleu_cercis} proves the BLEU--Cercis
isomorphism. Section~\ref{sec:multilingual} analyzes multilingual models
as multi-expert systems. Section~\ref{sec:codewitching} addresses
code-switching as gauge ambiguity. Section~\ref{sec:universal_grammar}
formulates UG as gauge-invariant content. Section~\ref{sec:evaluation_audit}
connects MT evaluation to the \SCX{} audit framework. Section~\ref{sec:embeddings}
analyzes word embeddings as gauge fields. Section~\ref{sec:nmt_gauge}
presents neural MT as learned gauge fixing. Section~\ref{sec:discussion}
discusses implications and open problems.

\begin{table}[H]
\centering
\caption{Complete dictionary: MT linguistics $\leftrightarrow$ SCX audit framework.}
\label{tab:mappings}
\begin{tabular}{@{}lll@{}}
\toprule
\textbf{\#} & \textbf{Linguistics / MT} & \textbf{SCX Audit} \\
\midrule
1 & Translation & Gauge transformation between coordinate charts \\
2 & BLEU score & \Cercis{} score ($Q$ = n-gram precision, $N$ = brevity penalty) \\
3 & mBERT/XLM-R shared latent space & Multi-expert shared representation \\
4 & Code-switching & Gauge ambiguity (undefined gauge at switch) \\
5 & Universal Grammar & Gauge-invariant linguistic substrate \\
6 & Human MT evaluation & $M$-expert audit with inter-annotator agreement \\
7 & Word embeddings & Gauge fields connecting language charts \\
8 & Neural MT attention & Learned gauge transformation \\
9 & Translationese & Gauge drift (incomplete gauge fixing) \\
\bottomrule
\end{tabular}
\end{table}


% ================================================================
\section{Translation as Gauge Transformation}
\label{sec:translation_gauge}
% ================================================================

\subsection{The Gauge Structure of Language}

In gauge theory, the fundamental objects are gauge-invariant---they do not
depend on the choice of local coordinates. The gauge field (connection)
mediates between coordinate charts, enabling parallel transport of
information. We apply this structure to language.

\begin{definition}[Translation as Gauge Transformation]
\label{def:translation_gauge}
Let $\Lang_s$ and $\Lang_t$ be source and target language coordinate
systems with meaning mappings $\Phi_s: \Vocab_s^* \to \M$ and
$\Phi_t: \Vocab_t^* \to \M$. A \textbf{translation} is a gauge
transformation $T_{s\to t}: \Sigma_s \to \Sigma_t$ such that for
all $x_s \in \Sigma_s$:
\begin{equation}
  \label{eq:translation_gauge}
  \Phi_t(T_{s\to t}(x_s)) = \Phi_s(x_s) + \delta_{\text{approx}},
\end{equation}
where $\delta_{\text{approx}} \in \M$ is an approximation error
(the information loss inherent in any finite translation). An
\textbf{exact translation} satisfies $\delta_{\text{approx}} = 0$.
\end{definition}

The gauge transformation $T_{s\to t}$ consists of a sequence of
local operations---word choice, syntactic restructuring, morphological
adjustment---each of which changes the surface form while preserving
the underlying meaning as closely as possible. This is exactly analogous
to a gauge transformation in physics: the content (the electromagnetic
field strength $F_{\mu\nu}$) is invariant, while the representation
(the vector potential $A_\mu$) changes.

\begin{definition}[Gauge-Invariant Content]
\label{def:gauge_invariant_content}
The \textbf{gauge-invariant content} of a sentence $x$ across a
set of languages $\{\Lang_a\}_{a\in A}$ is the equivalence class:
\begin{equation}
  \label{eq:gauge_invariant}
  \I(x) = \{m \in \M \mid \exists a \in A,\; \Phi_a^{-1}(m) \cap \Sigma_a \neq \emptyset\}
\end{equation}
under the equivalence relation $x_a \sim x_b \iff \Phi_a(x_a) = \Phi_b(x_b)$.
\end{definition}

\begin{remark}[No-Free-Lunch for Translation]
\label{rem:nfl_translation}
If $\mathcal{F}: \M \to \M$ is a nontrivial automorphism of meaning
space (a semantic transformation that preserves logical structure),
then translation $T_{s\to t}$ followed by $\mathcal{F}$ is also a
valid translation. This gauge freedom implies that translation quality
cannot be uniquely determined by meaning preservation alone---it must
additionally satisfy \emph{fluency} (gauge-fixing) constraints.
\end{remark}

\subsection{Parallel Transport of Meaning}

The gauge-theoretic interpretation provides a geometric picture of
translation as parallel transport along a path in ``language manifold'':

\begin{definition}[Language Manifold]
\label{def:language_manifold}
The \textbf{language manifold} $\LangManifold$ is a Riemannian manifold
whose points are language states and whose tangent space at a point
captures infinitesimal linguistic variation. A sentence in language
$L_a$ is a vector in $T_{L_a}\LangManifold$, and translation is
parallel transport along a curve $\gamma: [0,1] \to \LangManifold$
with $\gamma(0) = L_s$ and $\gamma(1) = L_t$.
\end{definition}

The parallel transport equation is:
\begin{equation}
  \label{eq:parallel_transport}
  \frac{D}{d\tau} T_{s\to t}(x_s) = 0,
\end{equation}
where $D/d\tau$ is the covariant derivative along $\gamma$ with
respect to a \textbf{linguistic connection} $\Gamma$ that encodes
cross-lingual correspondences. The connection $\Gamma$ is determined
by the bilingual training data: aligned sentence pairs define the
parallel transport map between language coordinate charts.

\begin{theorem}[Translation as Parallel Transport]
\label{thm:parallel_transport}
\rigorFull
For languages $L_s$ and $L_t$ with a sufficiently rich bilingual
corpus, the optimal translation function $T_{s\to t}^*$ is the
parallel transport map along the geodesic connecting $L_s$ and
$L_t$ in the language manifold $\LangManifold$, where the metric
is induced by the Fisher information of the bilingual distribution.
\end{theorem}

\begin{proof}
\noindent\textbf{Step 1: The bilingual distribution.}
Let $P(x_s, x_t)$ be the joint distribution over aligned sentence
pairs $(x_s, x_t) \in \Sigma_s \times \Sigma_t$. The conditional
distribution $P(x_t \mid x_s)$ defines the optimal stochastic
translation. The Fisher information metric on $\LangManifold$ at
language $L$ is:
\begin{equation}
  g_{ab}(L) = \E_{x \sim P_L}\left[
    \frac{\partial \log P_L(x)}{\partial \theta^a}
    \frac{\partial \log P_L(x)}{\partial \theta^b}
  \right],
\end{equation}
where $\theta^a$ are linguistic parameters (lexical, syntactic,
phonological).

\smallskip
\noindent\textbf{Step 2: Geodesic equation.}
The geodesic $\gamma(t)$ connecting $L_s$ to $L_t$ minimizes the
linguistic distance:
\begin{equation}
  d(L_s, L_t) = \int_0^1 \sqrt{g_{ab}(\gamma(t)) \dot\gamma^a(t) \dot\gamma^b(t)} \, dt.
\end{equation}
Parallel transport along this geodesic preserves the semantic
content up to second-order geodesic deviation.

\smallskip
\noindent\textbf{Step 3: Optimality.}
The optimal translation $T_{s\to t}^*$ is the solution to the
transport equation that minimizes the Kullback--Leibler divergence
between $P(x_t \mid x_s)$ and the parallel-transported distribution.
By the Amari--Chentsov theorem, the geodesic in Fisher metric is
the unique curve that minimizes this divergence under the Markov
condition $L_s \to T(L_s) \to L_t$.
\end{proof}

\subsection{Compositionality and Gauge Group}

Translations compose: translating from $L_s$ to $L_t$ and then to
$L_u$ should be equivalent to direct translation $L_s \to L_u$,
up to the approximation error. This is the group property of gauge
transformations:

\begin{proposition}[Translation Gauge Group]
\label{prop:gauge_group}
The set of translation gauge transformations $\{T_{a\to b}\}$ forms
a groupoid (a category with invertible morphisms) over the set of
languages. Composition satisfies:
\begin{equation}
  \label{eq:gauge_group}
  T_{b\to c} \circ T_{a\to b} \approx T_{a\to c},
\end{equation}
with equality only for exact translations. The deviation from
exact group composition is the \textbf{translation coherence gap}:
\begin{equation}
  \Delta_{\text{gap}}(a,b,c) = \Dist(T_{a\to c}, T_{b\to c} \circ T_{a\to b}),
\end{equation}
which measures how well the translation system respects the
groupoid structure.
\end{proposition}

\begin{remark}[Pivot Translation as Gauge Composition]
\label{rem:pivot}
In multilingual MT, translating through a pivot language $L_p$
(e.g., Korean $\to$ English $\to$ Vietnamese) exploits the
gauge groupoid structure. The quality loss incurred by pivot
translation is precisely the translation coherence gap
$\Delta_{\text{gap}}(a,p,b)$, which the \Cercis{} score can
measure as the drop in consensus between direct and pivot
translations.
\end{remark}


% ================================================================
\section{BLEU Score as Cercis Between Reference and Hypothesis}
\label{sec:bleu_cercis}
% ================================================================

\subsection{The BLEU Metric}

The Bilingual Evaluation Understudy (BLEU)~\cite{papineni2002bleu} is the
most widely used automatic metric for machine translation quality. It
computes the geometric mean of n-gram precisions $p_n$ (for $n = 1, 2, 3, 4$)
multiplied by a brevity penalty (BP) that penalizes hypotheses shorter than
the reference:

\begin{equation}
  \label{eq:bleu}
  \BLEU(r, h) = \text{BP}(r, h) \cdot \exp\left(\frac{1}{N}\sum_{n=1}^N \log p_n\right),
\end{equation}
where $p_n$ is the proportion of n-grams in the hypothesis $h$ that appear
in the reference $r$, and BP is the brevity penalty:
\begin{equation}
  \text{BP}(r, h) = \begin{cases}
    1 & \text{if } |h| \geq |r| \\
    \exp(1 - |r|/|h|) & \text{if } |h| < |r|.
  \end{cases}
\end{equation}

\subsection{The Cercis Score}

The \Cercis{} score in the \SCX{} framework is:
\begin{equation}
  \label{eq:cercis_general}
  S(s) = Q(s) + \eta(t) \cdot N(s),
\end{equation}
where $Q(s)$ is a quality (agreement) term and $N(s)$ is a novelty
(distinguishability) term modulated by an annealing schedule $\eta(t)$.

\begin{theorem}[BLEU = Cercis Score]
\label{thm:bleu_cercis}
\rigorFull
The BLEU score of a hypothesis translation $h$ against a reference
translation $r$ is a \Cercis{} score under the identifications:
\begin{equation}
  \label{eq:bleu_cercis_map}
  \begin{array}{ccl}
    Q_{\text{BLEU}}(r, h) &\longleftrightarrow& \displaystyle \frac{1}{N}\sum_{n=1}^N \log p_n \quad \text{(n-gram precision = quality)} \\[10pt]
    \eta(t) \cdot N_{\text{BLEU}}(r, h) &\longleftrightarrow& \log\text{BP}(r, h) \quad \text{(brevity penalty = novelty)} \\[6pt]
    \eta(t) &\longleftrightarrow& \mathbbm{1}_{|h| < |r|} \quad \text{(active only when hypothesis is shorter)}.
  \end{array}
\end{equation}
The total BLEU score in log-domain is $\log\BLEU(r, h) = Q_{\text{BLEU}} + \eta(t) \cdot N_{\text{BLEU}}$, which is structurally identical to
$\log S(s) = Q(s) + \eta(t) \cdot N(s)$.
\end{theorem}

\begin{proof}
\noindent\textbf{Step 1: Precision as quality.}
In the \SCX{} framework, the quality term $Q(s)$ measures agreement
between an expert's output and the consensus. In MT evaluation, the
reference translation $r$ is the ``expert consensus''---the gold
standard produced by a human translator. The hypothesis $h$ is the
``machine output.'' N-gram precision $p_n$ measures the proportion
of the machine's n-grams that agree with the reference:
\begin{equation}
  p_n = \frac{\sum_{\text{ngram} \in h} \mathbbm{1}_{\text{ngram} \in r}}{|h| - n + 1}.
\end{equation}
This is exactly a per-ngram quality score: $\log p_n$ is the
per-ngram log-agreement rate, analogous to the per-expert log-agreement
in the \Cercis{} score.

\smallskip
\noindent\textbf{Step 2: Brevity penalty as novelty.}
The brevity penalty activates when the hypothesis is shorter than
the reference ($|h| < |r|$). This penalizes the machine for
producing a ``novel'' (abbreviated) output that fails to convey
the full semantic content. In \Cercis{} terms:
\begin{equation}
  N_{\text{BLEU}} = \log\text{BP} = 1 - \frac{|r|}{|h|} \quad (< 0),
\end{equation}
which is a negative novelty score (a penalty for lacking expected
information). This is structurally identical to the novelty term
$\eta(t) \cdot N(s)$ in the \Cercis{} score when the annealing
schedule $\eta(t)$ is set to $\mathbbm{1}_{|h| < |r|}$---the
novelty term only contributes when the hypothesis is
``insufficiently expressive.''

\smallskip
\noindent\textbf{Step 3: Weighted geometric mean.}
The \Cercis{} score averages over $K$ experts. BLEU averages over
$N$ n-gram orders (usually $N = 4$). Both employ a geometric mean
structure: $\exp(\frac{1}{N}\sum \log p_n)$ is the geometric mean
of per-order precisions, exactly analogous to $\exp(Q(s) + \eta N(s))$
as the exponential of the mean log-score.

\smallskip
\noindent\textbf{Step 4: The reference as expert ensemble.}
The reference translation $r$ in BLEU plays the role of the expert
consensus in \Cercis{}. In practice, MT evaluation often uses
multiple references ($r_1, \ldots, r_M$) to capture the range of
acceptable translations. This is precisely an $M$-expert ensemble:
\begin{equation}
  \BLEU_M(r_1,\ldots,r_M, h) = \text{BP} \cdot \exp\left(\frac{1}{N}\sum_{n=1}^N \log p_n^{(M)}\right),
\end{equation}
where $p_n^{(M)}$ is the modified n-gram precision against the
multi-reference set (an n-gram counts as a match if it appears
in \emph{any} reference). This mirrors the multi-expert consensus
computation in \Cercis{}, where agreement with any subset of experts
can satisfy the quality condition.
\end{proof}

\begin{corollary}[Multi-Reference BLEU = Multi-Expert Cercis]
\label{cor:multiref_cercis}
When $M$ reference translations are available, the multi-reference
BLEU score is a \Cercis{} score computed over $M$ expert outputs,
with the consensus threshold adaptively set to the number of
references that contain each n-gram.
\end{corollary}

\subsection{Precision--Recall Decomposition}

The \Cercis{} score admits a decomposition into quality (precision-like)
and novelty (recall-like) components. In MT, this corresponds to
the precision--recall tradeoff:

\begin{definition}[N-Gram Cercis Decomposition]
\label{def:ngram_cercis}
For an MT system with hypothesis $h$ and reference $r$, define
the \textbf{n-gram Cercis score}:
\begin{equation}
  \label{eq:ngram_cercis}
  \Cercis_{\text{ngram}}(r, h) = \underbrace{\log\Precision_{\text{ngram}}(h, r)}_{Q(s)}
    + \eta(t) \cdot \underbrace{\log\Recall_{\text{ngram}}(h, r)}_{N(s)},
\end{equation}
where:
\begin{equation}
  \Precision_{\text{ngram}}(h, r) = \frac{|h \cap r|}{|h|},
  \qquad
  \Recall_{\text{ngram}}(h, r) = \frac{|h \cap r|}{|r|},
\end{equation}
and $|h \cap r|$ is the count of n-grams appearing in both hypothesis
and reference, $|h|$ is the total n-gram count in the hypothesis,
and $|r|$ is the total in the reference.
\end{definition}

\begin{proposition}[BLEU = Precision-Dominated Cercis]
\label{prop:bleu_precision_cercis}
Standard BLEU is a precision-dominated \Cercis{} score: it weights
precision ($Q$) more strongly than recall ($N$). The brevity penalty
$\log\text{BP}$ provides the only recall-based correction. Other
metrics (e.g., ROUGE~\cite{lin2004rouge}, METEOR~\cite{banerjee2005meteor},
COMET~\cite{rei2020comet}) correspond to different quality--novelty
tradeoffs in the \Cercis{} framework:
\begin{equation}
  \begin{array}{ccl}
    \text{BLEU} &:& \eta(t) = \mathbbm{1}_{|h| < |r|} \quad \text{(recall only when short)} \\[4pt]
    \text{ROUGE-L} &:& \eta(t) = 1 \quad \text{(balanced precision--recall)} \\[4pt]
    \text{METEOR} &:& \eta(t) = \frac{9\cdot\Precision \cdot \Recall}{\Precision + 9\Recall} \quad \text{(F-measure weighted toward recall)}
  \end{array}
\end{equation}
\end{proposition}


% ================================================================
\section{Multilingual Models as Multi-Expert Shared Latent Space}
\label{sec:multilingual}
% ================================================================

Multilingual pre-trained language models such as mBERT~\cite{devlin2019bert}
and XLM-R~\cite{conneau2020unsupervised} learn a shared representation
space across dozens of languages. We show that this shared representation
is precisely the \SCX{} multi-expert consensus space.

\subsection{The Architecture of Multilingual Models}

\begin{definition}[Multilingual Encoder as Multi-Expert System]
\label{def:multilingual_encoder}
A multilingual language model with $N$ languages and $K$ transformer
layers defines a set of \textbf{language-specific experts}
$\{E_a\}_{a=1}^N$ where each expert $E_a$ is a function:
\begin{equation}
  E_a(x_a) = \text{FFN}_{\text{shared}} \circ \text{SelfAttn}_{\text{shared}} \circ \text{Embed}_a(x_a),
\end{equation}
where $\text{Embed}_a$ is a language-specific embedding layer (gauge
chart), and the shared transformer layers (SelfAttn, FFN) encode
gauge-invariant features.
\end{definition}

\begin{theorem}[Cross-Lingual Transfer = Gauge-Invariant Feature Extraction]
\label{thm:crosslingual_transfer}
\rigorFull
For languages $L_s$ and $L_t$, a multilingual model achieves
zero-shot cross-lingual transfer if and only if the shared encoder
$\psi: \Vocab^* \to \R^d$ extracts features that are
\textbf{gauge-invariant}---i.e., for any sentence $x_s \in \Sigma_s$
and its translation $x_t \in \Sigma_t$:
\begin{equation}
  \label{eq:crosslingual_invariant}
  \|\psi(x_s) - \psi(x_t)\| \leq \varepsilon_{\text{transfer}},
\end{equation}
where $\varepsilon_{\text{transfer}}$ is the cross-lingual transfer
gap. The shared representation space $\R^d$ is the \textbf{gauge-invariant
latent space} that approximates the meaning manifold $\M$.
\end{theorem}

\begin{proof}
\noindent\textbf{Step 1: Shared vocabulary of subwords.}
Multilingual models tokenize all languages using a shared subword
vocabulary (e.g., SentencePiece~\cite{kudo2018sentencepiece}). This
creates a common token-level representation: the embedding lookup
$E: \Vocab_{\text{shared}} \to \R^d$ maps subword tokens to vectors
in a shared space. Different languages map to overlapping regions of
this shared token space, ensuring that the embedding layer is
approximately gauge-invariant.

\smallskip
\noindent\textbf{Step 2: Language-agnostic attention.}
The self-attention mechanism operates on token representations
without language identity information (after the initial embedding
and position encoding). The attention pattern between tokens is
predominantly determined by syntactic and semantic relationships
that are shared across languages---subject--verb agreement, modifier--head
relations, coreference---all of which are gauge-invariant features.

\smallskip
\noindent\textbf{Step 3: Contrastive alignment.}
The model is trained on masked language modeling (MLM) across
concatenated corpora from all languages. This forces the model to
predict masked tokens using cross-lingual context. The MLM objective:
\begin{equation}
  \loss_{\text{MLM}} = -\E_{x} \sum_{i \in \text{masked}}
    \log P(x_i \mid x_{\backslash i}),
\end{equation}
implicitly aligns the representation spaces: the same semantic role
in different languages maps to similar hidden representations because
the model must use context from one language to predict tokens in
another. This alignment is precisely the parallel transport condition
from Section~\ref{sec:translation_gauge}.

\smallskip
\noindent\textbf{Step 4: The transfer gap.}
The inequality $\|\psi(x_s) - \psi(x_t)\| \leq \varepsilon_{\text{transfer}}$
holds because the shared encoder $\psi$ is a \textbf{contraction}
on the meaning manifold: it maps semantically equivalent sentences
from different languages to nearby points in the shared representation
space. The magnitude of $\varepsilon_{\text{transfer}}$ is bounded by
the degree of language separation in the training data and the
dimensionality $d$ of the shared space:
\begin{equation}
  \varepsilon_{\text{transfer}} \leq \frac{C}{\sqrt{d}} \cdot \sum_{a=1}^N
    \text{KL}(P_a \| P_{\text{avg}}),
\end{equation}
where $\text{KL}(P_a \| P_{\text{avg}})$ is the Kullback--Leibler
divergence between the monolingual distribution of language $L_a$
and the average distribution across all languages.
\end{proof}

\subsection{The Approximate UG Manifold}

\begin{proposition}[Multilingual Model Latent Space = Approximate UG Manifold]
\label{prop:ug_manifold}
The shared latent space of a multilingual model (e.g., the
[CLS] token representation from the final layer of mBERT)
forms an \textbf{approximate Universal Grammar manifold}:
it encodes linguistic features that are invariant across all
languages in its training set, up to the resolution limit of
the model. For $N$ languages, the latent space $\R^d$ is a
low-dimensional projection of the $N$-fold product of
language-specific feature spaces, with the projection
minimizing the cross-lingual reconstruction loss:
\begin{equation}
  \min_{\psi, \{\phi_a\}} \sum_{a=1}^N \E_{x_a \sim P_a}
    \|x_a - \phi_a(\psi(x_a))\|^2,
\end{equation}
where $\phi_a: \R^d \to \Vocab_a^*$ is the language-specific
decoder (inverse gauge transformation).
\end{proposition}

\begin{proof}[Proof Sketch]
\noindent\textbf{Step 1: Autoencoder structure.}
The multilingual MLM objective defines an implicit autoencoder:
$\psi$ encodes to the shared space, and a language-specific head
decodes to the original language. This is exactly the gauge-theoretic
structure: $\psi$ extracts gauge-invariant features, and the
language-specific heads implement the gauge fixing.

\smallskip
\noindent\textbf{Step 2: Representation topology.}
Neural representations of different languages in the shared space
form partially overlapping clusters. The overlap regions correspond
to linguistic features shared across languages (universals). The
Hausdorff distance between language clusters in the shared space
measures the linguistic distance between languages:
\begin{equation}
  d_H(L_a, L_b) = \max\left\{
    \sup_{v \in \psi(\Sigma_a)} \inf_{w \in \psi(\Sigma_b)} \|v - w\|,\;
    \sup_{w \in \psi(\Sigma_b)} \inf_{v \in \psi(\Sigma_a)} \|v - w\|
  \right\}.
\end{equation}
\end{proof}


% ================================================================
\section{Code-Switching as Gauge Ambiguity}
\label{sec:codewitching}
% ================================================================

\subsection{The Phenomenon of Code-Switching}

Code-switching (CS) is the alternation between two or more languages
within a single utterance, conversation, or text~\cite{myers2002code}.
It is prevalent in multilingual communities worldwide: Spanish--English
in the United States, Hindi--English in India, Mandarin--Cantonese in
Southeast Asia, and dozens of other language pairs.

\subsection{Gauge Ambiguity at Code-Switch Points}

\begin{definition}[Code-Switching as Gauge Ambiguity]
\label{def:codewitching_gauge}
Let $L_a$ and $L_b$ be two languages with coordinate charts
$(\Vocab_a, \Phi_a)$ and $(\Vocab_b, \Phi_b)$. A \textbf{code-switched
utterance} $x_{ab}$ is a sequence:
\begin{equation}
  x_{ab} = x_a^{(1)} \oplus x_b^{(1)} \oplus x_a^{(2)} \oplus x_b^{(2)} \oplus \cdots,
\end{equation}
where $x_a^{(i)} \in \Sigma_a$, $x_b^{(j)} \in \Sigma_b$, and
$\oplus$ denotes concatenation. At each switch point $\oplus$,
the gauge is \textbf{undefined}: the meaning mapping $\Phi$
momentarily lacks a well-defined coordinate chart.
\end{definition}

\begin{theorem}[Code-Switching = Gauge Singularity]
\label{thm:codewitching_singularity}
\rigorFull
At a code-switch point between languages $L_a$ and $L_b$, the
gauge transformation $T_{a\to b}$ is undefined if and only if
the languages differ in their syntactic parameter values at
the switch boundary. The degree of gauge ambiguity is measured
by the \textbf{CS perplexity}:
\begin{equation}
  \label{eq:cs_perplexity}
  \Pbb_{\text{CS}}(x_{ab}) = \exp\left(-\sum_{i=1}^{|x_{ab}|}
    \log P_{\text{LM}}(x_{ab}^{(i)} \mid x_{ab}^{(<i)})\right),
\end{equation}
where $P_{\text{LM}}$ is a language model trained on monolingual
data only. The CS perplexity is strictly higher than the perplexity
of equivalent monolingual utterances in either language:
\begin{equation}
  \Pbb_{\text{CS}}(x_{ab}) > \max\{\Pbb_{L_a}(x_a), \Pbb_{L_b}(x_b)\},
\end{equation}
for any $x_a, x_b$ that express the same content as $x_{ab}$.
\end{theorem}

\begin{proof}
\noindent\textbf{Step 1: Syntactic parameter mismatch.}
Languages differ in syntactic parameters such as head direction
(head-initial vs.\ head-final), pro-drop, and word order (SVO vs.\ SOV).
At a switch point, the grammar of language $L_a$ may require a
different syntactic structure than the grammar of $L_b$. For example,
in Spanish--English code-switching:
\begin{equation}
  \text{``I want } \underbrace{\text{comprar}}_{V_{\text{SP}}} \text{ the book''}
\end{equation}
has a mixed VP structure: the verb is Spanish (head-initial) but the
object NP follows English word order. This syntactic mismatch means
no single set of grammatical rules can parse the utterance---the
gauge transformation $T_{\text{EN}\to\text{ES}}$ is undefined at
the switch.

\smallskip
\noindent\textbf{Step 2: The Matrix Language Framework.}
Myers-Scotton's Matrix Language Framework~\cite{myers2002code} posits
that CS utterances have a \textbf{matrix language} (the dominant grammar)
and an \textbf{embedded language} (lexical insertions). The matrix
language provides the syntactic frame, and the embedded language
provides content morphemes. In gauge-theoretic terms, the matrix
language fixes the gauge, and embedded language elements are
gauge-covariant insertions: they transform under the matrix language's
syntactic rules.

\smallskip
\noindent\textbf{Step 3: Perplexity bound.}
A monolingual LM trained on $L_a$ will assign near-zero probability
to $L_b$ tokens following $L_a$ syntactic frames, and vice versa.
The CS perplexity $\Pbb_{\text{CS}}(x_{ab})$ is therefore bounded
from below by the probability of the switch conditioned on the
bilingual context, which is always strictly positive but less than
the best monolingual probability.
\end{proof}

\begin{corollary}[CS Detection = Gauge Ambiguity Detection]
\label{cor:cs_detection}
Code-switched utterances can be detected by measuring the \textbf{gauge
ambiguity score}:
\begin{equation}
  \label{eq:gauge_ambiguity}
  G(x) = \frac{\max_{a,b} \Pbb_{\text{CS}}(x_{ab})}{\Pbb_{\text{mono}}(x)},
\end{equation}
where $\Pbb_{\text{mono}}(x)$ is the maximum probability assigned
to $x$ by any monolingual model. If $G(x) > \tau_{\text{CS}}$, the
utterance contains a code switch at confidence $1-\delta$.
\end{corollary}

\subsection{Code-Switching Cercis Score}

In the \SCX{} audit framework, code-switching can be evaluated by
measuring the \Cercis{} score between the CS utterance and its
monolingual paraphrases:

\begin{proposition}[CS Cercis Score]
\label{prop:cs_cercis}
For a code-switched utterance $x_{ab}$, the \Cercis{} score against
monolingual paraphrases $x_a$ and $x_b$ is:
\begin{equation}
  S_{\text{CS}}(x_{ab}) = \frac{1}{2}\left[S(x_{ab}, x_a) + S(x_{ab}, x_b)\right],
\end{equation}
where $S(x_{ab}, x_a)$ is the \Cercis{} score (e.g., BLEU) between
the CS utterance and its monolingual counterpart. A low $S_{\text{CS}}$
indicates that the CS utterance is not a faithful expression of the
same semantic content in either language---it occupies a ``gauge
singularity'' between the two coordinate charts.
\end{proposition}


% ================================================================
\section{Universal Grammar as the Gauge-Invariant Substrate}
\label{sec:universal_grammar}
% ================================================================

\subsection{Principles and Parameters as Gauge Degrees of Freedom}

Chomsky's theory of Universal Grammar~\cite{chomsky1965aspects,
chomsky1981lectures} posits that all human languages share a common
set of principles (invariant constraints) while differing in a finite
set of parameters (binary switches that determine surface syntactic
variation).

\begin{definition}[Universal Grammar as Gauge-Invariant Content]
\label{def:ug_gauge}
Let $\{\Lang_a\}_{a=1}^N$ be the set of all human languages, each
a coordinate chart on the meaning manifold $\M$. The \textbf{Universal
Grammar} $\UG$ is the intersection of gauge-invariant linguistic
features across all language charts:
\begin{equation}
  \label{eq:ug_invariant}
  \UG = \bigcap_{a=1}^N \left\{ f: \Sigma_a \to \R \;\big|\;
    f \text{ is invariant under } T_{a\to b}\; \forall b \right\}.
\end{equation}
That is, $\UG$ is the set of features $f$ such that for any sentence
$x_a \in \Sigma_a$ and its translation $x_b = T_{a\to b}(x_a)$,
$f(x_a) = f(x_b)$. These are the \textbf{absolute linguistic invariants}.
\end{definition}

\begin{theorem}[Universal Grammar = Gauge-Invariant Substrate]
\label{thm:ug_invariant}
\rigorFull
The Universal Grammar $\UG$ is the \textbf{gauge-invariant substrate}
of language: the maximal set of linguistic properties that remain
unchanged under all possible language-to-language gauge transformations.
The principles of UG are gauge constraints, and the parameters are
gauge degrees of freedom:
\begin{equation}
  \label{eq:principles_parameters}
  \UG = \underbrace{\{\text{principles}\}}_{\text{gauge constraints}} \cup
  \underbrace{\{\text{parameter choices}\}}_{\text{gauge-fixing conditions}}.
\end{equation}
\end{theorem}

\begin{proof}
\noindent\textbf{Step 1: Principles as gauge constraints.}
Chomsky's principles---structure-dependency, binding theory, X-bar
schema, subjacency---are constraints that hold in \emph{every} language.
If a principle $P$ held in $L_a$ but not in $L_b$, then there would
exist a sentence $x_a \in \Sigma_a$ satisfying $P$ whose translation
$x_b = T_{a\to b}(x_a)$ violates $P$. But then $T_{a\to b}$ would
not be a valid gauge transformation (it would change the gauge-invariant
content), contradicting Definition~\ref{def:translation_gauge}.
Thus principles are gauge-invariant by construction.

\smallskip
\noindent\textbf{Step 2: Parameters as gauge degrees of freedom.}
Parameters---such as the null subject parameter (pro-drop:
Italian/ Spanish vs.\ English/French), the head direction parameter
(head-initial English vs.\ head-final Japanese), and the wh-movement
parameter (overt movement in English vs.\ in-situ in Chinese)---are
the \textbf{gauge degrees of freedom} that distinguish language
coordinate charts. Two languages differing in a parameter value
are using different gauge fixings of the same underlying principle.

\smallskip
\noindent\textbf{Step 3: The invariance condition.}
Formally, if $\theta_a = \{\pi_a^{(1)}, \ldots, \pi_a^{(k)}\}$ are
the parameter values of language $L_a$, then the gauge transformation
$T_{a\to b}$ must satisfy:
\begin{equation}
  T_{a\to b}(\theta_a) = \theta_b,
\end{equation}
meaning that the translation function $T_{a\to b}$ maps between
parameter configurations. The UG is the set of constraints that
do not depend on $\theta$:
\begin{equation}
  \UG = \{P \mid \forall \theta, P(\theta) = \text{true}\}.
\end{equation}
\end{proof}

\begin{corollary}[Learnability from Gauge Invariance]
\label{cor:learnability}
The gauge invariance of UG implies that language acquisition can be
modeled as a \textbf{gauge-fixing process}: the child starts with
the gauge-invariant principles (UG) and infers the parameter values
(grammar) from exposure to a particular language. The number of
parameters $k$ determines the VC dimension of the hypothesis space,
with $|\Theta| = 2^k$ possible grammars.
\end{corollary}

\subsection{Empirical Predictions}

The gauge-theoretic formulation of UG makes testable predictions:

\begin{enumerate}[leftmargin=*, label=(\arabic*)]
  \item \textbf{Transfer asymmetry}: Cross-lingual transfer should be
        easier between languages that share more parameter values
        (gauge fixings are closer), consistent with empirical findings
        in multilingual NLP~\cite{pires2019multilingual,wu2024bilingual}.
  \item \textbf{Parameter clustering}: Languages should cluster in
        parameter space according to their phylogenetic families,
        because languages within a family share inherited parameter
        values (gauge fixings). This matches the observed typological
        clusters in the World Atlas of Language Structures
        (WALS)~\cite{dryer2013wals}.
  \item \textbf{Impossible languages}: Not all combinations of
        parameter values are realizable, because the parameter
        space is constrained by the principles (gauge constraints).
        This matches Greenberg's linguistic universals~\cite{greenberg1963universals}.
\end{enumerate}


% ================================================================
\section{Machine Translation Evaluation as Audit}
\label{sec:evaluation_audit}
% ================================================================

\subsection{Human Evaluation as Multi-Expert Audit}

Human evaluation of MT quality---using frameworks such as MQM
(Multidimensional Quality Metrics)~\cite{lommel2014multidimensional}
and DQF (Dynamic Quality Framework)~\cite{castilho2017guide}---is
a direct instance of the \SCX{} audit protocol.

\begin{definition}[MT Evaluation as Audit]
\label{def:mt_audit}
An MT evaluation with $M$ human annotators and $K$ quality dimensions
is a \textbf{multi-expert audit} where:
\begin{itemize}[leftmargin=*]
  \item Each annotator $a_m$ is an \textbf{expert} whose gatekeeper
        parameter $g_m$ encodes their linguistic background, preferences,
        and biases;
  \item Each quality dimension $d_k$ (e.g., adequacy, fluency,
        terminology, style) is a \textbf{probe direction};
  \item The aggregation of annotator scores is the \textbf{consensus
        Cercis score};
  \item Annotator disagreement is a measure of \textbf{bias radiation}
        (Section~\ref{sec:bleu_cercis}).
\end{itemize}
\end{definition}

\begin{theorem}[Inter-Annotator Agreement = Cercis Score]
\label{thm:inter_annotator}
\rigorFull
For an MT evaluation with $M$ annotators rating $N$ translations
on a scale $\{1, \ldots, L\}$, the inter-annotator agreement
coefficient $\kappa$ (e.g., Fleiss' kappa) is a \Cercis{} score:
\begin{equation}
  \label{eq:kappa_cercis}
  \kappa = \frac{\Cercis_{\text{annotators}} - \E[\Cercis_{\text{annotators}} \mid H_0]}
    {1 - \E[\Cercis_{\text{annotators}} \mid H_0]},
\end{equation}
where $\Cercis_{\text{annotators}}$ is the average pairwise agreement
rate, and $H_0$ is the null hypothesis of random annotation. The
quality term $Q$ is the observed agreement, and the novelty term
$N$ captures systematic disagreement patterns (bias).
\end{theorem}

\begin{proof}
\noindent\textbf{Step 1: Fleiss' kappa as Cercis.}
Fleiss' kappa for $M$ annotators on $N$ items with $L$ categories is:
\begin{equation}
  \kappa = \frac{\bar{P} - \bar{P}_e}{1 - \bar{P}_e},
\end{equation}
where $\bar{P}$ is the observed agreement and $\bar{P}_e$ is the
expected agreement under the null (random annotation). In the
\Cercis{} framework:
\begin{equation}
  \bar{P} = \frac{1}{N} \sum_{i=1}^N \frac{1}{M(M-1)} \sum_{j=1}^M
    \sum_{k \neq j} \mathbbm{1}[a_{ij} = a_{ik}],
\end{equation}
is the quality score $Q$---the average pairwise agreement. The
denominator $1 - \bar{P}_e$ normalizes by the maximum possible
disagreement.

\smallskip
\noindent\textbf{Step 2: Bias detection.}
When $\kappa$ is low relative to a threshold $\kappa_0$, the audit
flags the evaluation as unreliable---annotator bias is high, and
the \Cercis{} score is dominated by the novelty term $N$ (the
systematic disagreement pattern). This is exactly the \SCX{} bias
detection mechanism: when expert outputs diverge beyond the expected
rate, the gatekeeper parameter $g$ is inferred to be large.

\smallskip
\noindent\textbf{Step 3: MQM as multi-probe audit.}
MQM evaluates translations across multiple error categories
(accuracy, fluency, terminology, style, locale convention).
Each category is an independent audit probe. The total audit
score is the weighted \Cercis{} across categories:
\begin{equation}
  S_{\text{MQM}} = \sum_{k=1}^K w_k \cdot \Cercis_k(\text{annotators}),
\end{equation}
where $w_k$ are category weights and $\Cercis_k$ is the per-category
agreement score.
\end{proof}

\subsection{Automatic Evaluation as Audit}

In the \SCX{} framework, automatic MT metrics (BLEU, TER, chrF,
COMET) are \textbf{audit instruments}---scoring functions that
probe the hypothesis translation for agreement with the reference
consensus. Table~\ref{tab:metrics} maps each metric to its \Cercis{}
interpretation.

\begin{table}[H]
\centering
\caption{Automatic MT metrics as \Cercis{} audit instruments.}
\label{tab:metrics}
\begin{tabular}{@{}llll@{}}
\toprule
\textbf{Metric} & \textbf{Quality ($Q$)} & \textbf{Novelty ($N$)} & \textbf{Gauge Interpretation} \\
\midrule
BLEU & N-gram precision & Brevity penalty & Precision-biased Cercis \\
chrF & Character n-gram F-score & (balanced F) & Character-level Cercis \\
TER & Edit distance \emph{inverse} & (none) & Purely quality-based \\
METEOR & Unigram precision & Recall + fragment penalty & Balanced Cercis \\
COMET & Neural similarity & (learned) & Learned Cercis on latent space \\
\bottomrule
\end{tabular}
\end{table}

\subsection{The Audit Protocol for MT Evaluation}

We formalize the optimal MT evaluation protocol using the \SCX{}
audit framework:

\begin{protocol}[Optimal MT Audit]
\label{proto:mt_audit}
Given a hypothesis translation $h$ and a set of references $R = \{r_1, \ldots, r_M\}$:
\begin{enumerate}[leftmargin=*, label=\textbf{Step \arabic*:}]
  \item \textbf{Expert ensemble formation}: Collect $M$ human reference
        translations and/or automatic metric scores.
  \item \textbf{Cercis score computation}: Compute the multi-reference
        \Cercis{} score $S_{\text{MT}}(h, R) = Q(h, R) + \eta(t) N(h, R)$.
  \item \textbf{Confidence assessment}: Compute the audit confidence
        interval $\CI_{\delta}(S_{\text{MT}})$ using the Hoeffding bound
        on annotation variance.
  \item \textbf{Bias diagnosis}: If the inter-annotator agreement
        $\kappa_{\text{annotator}} < \kappa_{\text{threshold}}$, flag
        potential annotator bias and compute the bias-corrected
        \Cercis{} score $\tilde{S}_{\text{MT}}$.
  \item \textbf{Gauge checking}: Verify that the \Cercis{} score is
        approximately invariant under the choice of reference language
        (if evaluating in a multilingual setting): for any $r_a, r_b \in R$
        from different languages, $S(h, r_a) \approx S(h, r_b)$.
\end{enumerate}
\end{protocol}


% ================================================================
\section{Word Embeddings as Gauge Fields}
\label{sec:embeddings}
% ================================================================

\subsection{Monolingual Embeddings as Local Gauge Fields}

Word embeddings (word2vec~\cite{mikolov2013distributed}, GloVe~\cite{pennington2014glove},
FastText~\cite{bojanowski2017enriching}) map each word in a vocabulary
to a vector in $\R^d$. In gauge-theoretic language, each embedding
space is a \textbf{local gauge field}: a fiber bundle over the lexical
inventory, where the fiber at each word is its semantic vector.

\begin{definition}[Word Embedding as Gauge Field]
\label{def:embedding_gauge}
Let $\Vocab_a$ be the vocabulary of language $L_a$. A \textbf{word
embedding} is a gauge field $A_a: \Vocab_a \to \R^d$ that assigns
to each word $w \in \Vocab_a$ a vector $A_a(w) \in \R^d$ such that
the inner product $\langle A_a(w_1), A_a(w_2) \rangle$ measures
the semantic similarity between $w_1$ and $w_2$ in the coordinate
chart of $L_a$. The gauge field transforms under language
translation as:
\begin{equation}
  \label{eq:embedding_transform}
  A_b(T_{a\to b}(w)) = U_{b\leftarrow a} \cdot A_a(w) + \delta_b,
\end{equation}
where $U_{b\leftarrow a} \in O(d)$ is a rotation matrix (the
\textbf{parallel transport operator}) and $\delta_b \in \R^d$
is a translation vector (the \textbf{coordinate offset}).
\end{definition}

\subsection{Cross-Lingual Embedding Alignment as Gauge Fixing}

\begin{definition}[Cross-Lingual Alignment as Gauge Connection]
\label{def:alignment_gauge}
Let $\Embed_a: \Vocab_a \to \R^d$ and $\Embed_b: \Vocab_b \to \R^d$
be embedding gauge fields for languages $L_a$ and $L_b$. A
\textbf{cross-lingual alignment} is a gauge connection:
\begin{equation}
  \Align_{a\to b}: \R^d \to \R^d,
\end{equation}
typically a linear map $\Align_{a\to b}(v) = W_{a\to b} v + b_{a\to b}$
learned from a bilingual dictionary. The alignment is a \textbf{gauge
fixing} if it minimizes the semantic distortion:
\begin{equation}
  \label{eq:alignment_gauge_fixing}
  W_{a\to b}^* = \argmin_{W \in \R^{d\times d}} \sum_{(w_a, w_b) \in \D_{\text{biling}}}
    \|W \Embed_a(w_a) - \Embed_b(w_b)\|^2.
\end{equation}
\end{definition}

\begin{theorem}[Embedding Alignment = Gauge Connection Learning]
\label{thm:embedding_alignment}
\rigorFull
Cross-lingual word embedding alignment via linear mapping (e.g.,
MUSE~\cite{conneau2018word}, VecMap~\cite{artetxe2018massively})
learns a gauge connection between language coordinate charts.
The quality of the alignment---measured by the bilingual dictionary
induction score (P@k)---is a \Cercis{} score between the aligned
and gold-standard bilingual lexicons.
\end{theorem}

\begin{proof}
\noindent\textbf{Step 1: Procrustes alignment.}
The optimal alignment $W_{a\to b}$ under the Frobenius norm constraint
$W^T W = I$ is solved by the orthogonal Procrustes problem:
\begin{equation}
  W_{a\to b}^* = UV^T, \quad \text{where } U \Sigma V^T = \text{SVD}(X_a X_b^T),
\end{equation}
with $X_a \in \R^{k \times d}$ and $X_b \in \R^{k \times d}$ being matrices
of aligned word embeddings. This solution is the \textbf{parallel
transport operator} $U_{b\leftarrow a}$ from Definition~\ref{def:embedding_gauge}.

\smallskip
\noindent\textbf{Step 2: P@k as Cercis.}
The precision@k for bilingual lexicon induction is:
\begin{equation}
  \text{P@k} = \frac{1}{k} \sum_{i=1}^k \mathbbm{1}[y_i \in Y_{\text{gold}}(x_i)],
\end{equation}
where $Y_{\text{gold}}(x_i)$ is the set of gold translations for word $x_i$.
This is exactly a quality score $Q$ in the \Cercis{} framework: it measures
the proportion of aligned pairs that match the expert consensus (the gold
lexicon). The \Cercis{} score for alignment is:
\begin{equation}
  S_{\text{align}} = Q_{\text{P@k}} + \eta(t) \cdot N_{\text{diversity}},
\end{equation}
where $N_{\text{diversity}}$ penalizes trivial alignments (e.g., all words
aligned to the same target word).

\smallskip
\noindent\textbf{Step 3: Invariance under gauge transformation.}
The alignment $W_{a\to b}$ satisfies gauge covariance:
\begin{equation}
  W_{b\to c} \circ W_{a\to b} \approx W_{a\to c},
\end{equation}
up to the alignment coherence gap (Proposition~\ref{prop:gauge_group}).
This ensures that the aligned embedding space is approximately
gauge-invariant: semantic relations are preserved under composition
of alignments.
\end{proof}

\begin{corollary}[Language Distance from Alignment Error]
\label{cor:align_distance}
The Frobenius norm of the residual alignment error:
\begin{equation}
  d_{\text{align}}(L_a, L_b) = \min_{W \in O(d)} \|W X_a - X_b\|_F,
\end{equation}
defines a metric on the language manifold $\LangManifold$. This metric
correlates with phylogenetic language distance and with the BLEU score
between language pairs.
\end{corollary}


% ================================================================
\section{Neural Machine Translation as Learned Gauge Fixing}
\label{sec:nmt_gauge}
% ================================================================

\subsection{The Transformer as Gauge-Equivariant Architecture}

The transformer attention mechanism~\cite{vaswani2017attention} is
naturally gauge-equivariant: its outputs transform covariantly under
input reparameterizations.

\begin{definition}[Attention as Gauge Transformation]
\label{def:attention_gauge}
The scaled dot-product attention computes:
\begin{equation}
  \text{Attention}(Q, K, V) = \softmax\left(\frac{QK^T}{\sqrt{d_k}}\right) V,
\end{equation}
where $Q = X W_Q$, $K = X W_K$, $V = X W_V$ are linear projections
of the input sequence $X \in \R^{n \times d}$. If the input transforms
under a gauge transformation $X \to X' = X \cdot G$ (where $G \in GL(d)$
is an invertible matrix), then:
\begin{equation}
  \text{Attention}(Q', K', V') = \text{Attention}(Q, K, V) \cdot G_V,
\end{equation}
where $G_V$ is the projection of $G$ onto the value subspace. The
attention mechanism is \textbf{gauge-covariant}: the learned alignment
patterns are invariant under gauge transformations of the input space.
\end{definition}

\begin{theorem}[Transformer = Gauge-Equivariant Architecture]
\label{thm:transformer_gauge}
\rigorFull
The transformer encoder--decoder architecture is \textbf{gauge-equivariant}:
for any input sequence $x_s$ in source language $L_s$, and any gauge
transformation $G$ of the shared embedding space, the encoder output
$\psi(x_s)$ and decoder output $T_{s\to t}(x_s)$ satisfy:
\begin{equation}
  \label{eq:transformer_gauge_eq}
  \psi(G \cdot x_s) = G_{\psi} \cdot \psi(x_s), \qquad
  T_{s\to t}(G \cdot x_s) = G_T \cdot T_{s\to t}(x_s),
\end{equation}
where $G_{\psi}$ and $G_T$ are induced representations of $G$ in
the encoder and decoder feature spaces, respectively. This
equivariance ensures that the translation quality is independent
of the specific gauge (language) chosen for representation.
\end{theorem}

\begin{proof}[Proof Sketch]
\noindent\textbf{Step 1: Layer normalization as gauge-fixing.}
Layer normalization~\cite{ba2016layer} computes:
\begin{equation}
  \text{LayerNorm}(x) = \frac{x - \mu}{\sqrt{\sigma^2 + \epsilon}} \odot \gamma + \beta.
\end{equation}
This removes the global gauge (mean and variance) of the representation,
ensuring that gauge transformations preserving the mean and variance
structure are factored out. LayerNorm is a \textbf{gauge-fixing
operation}.

\smallskip
\noindent\textbf{Step 2: Residual connections as gauge invariance.}
The residual connection $x_{l+1} = x_l + F(x_l)$ ensures that the
representation space at each layer is a \textbf{trivial bundle}:
the fiber at position $l$ is identified with the fiber at position
$l+1$ via the identity map. This ensures that gauge transformations
are consistently transported across layers.

\smallskip
\noindent\textbf{Step 3: Position encoding as coordinate chart.}
Sinusoidal position encodings provide a global coordinate chart on
the sequence manifold. The encoding is gauge-covariant: shifting
positions by $k$ corresponds to rotating the encoding by a fixed
angle, which is a gauge transformation of the position space.
\end{proof}

\subsection{Attention as Learned Gauge Transformation}

In the encoder--decoder attention, the decoder queries attend to
encoder keys over the source sequence. This cross-attention is the
learned gauge transformation:

\begin{proposition}[Cross-Attention = Gauge Transformation]
\label{prop:cross_attention}
The encoder--decoder cross-attention in a transformer NMT system
computes the gauge transformation $T_{s\to t}$ at each decoding
step. The attention weights $a_{ij}$ define the \textbf{gauge
connection coefficients}:
\begin{equation}
  a_{ij} = \frac{\exp(e_{ij})}{\sum_k \exp(e_{ik})}, \quad
  e_{ij} = \frac{(h_i W_Q)(s_j W_K)^T}{\sqrt{d_k}},
\end{equation}
where $h_i$ is the $i$-th decoder hidden state and $s_j$ is the
$j$-th encoder hidden state. The attention-weighted sum:
\begin{equation}
  c_i = \sum_j a_{ij} (s_j W_V)
\end{equation}
is the parallel transport of source information to the target side,
with $a_{ij}$ as the transport coefficients.
\end{proposition}

\subsection{Backtranslation and Iterative Gauge Refinement}

Backtranslation~\cite{sennrich2016improving}---translating target-side
monolingual data back to the source side---is an iterative gauge
refinement process:

\begin{proposition}[Backtranslation as Gauge Self-Consistency]
\label{prop:backtranslation}
Backtranslation enforces the gauge-fixing condition:
\begin{equation}
  T_{t\to s}(T_{s\to t}(x_s)) \approx x_s,
\end{equation}
which is the \textbf{invertibility constraint} on the gauge transformation.
The reconstruction error:
\begin{equation}
  \loss_{\text{BT}} = \E_{x_s \sim P_s}[\Dist(x_s, T_{t\to s}(T_{s\to t}(x_s)))],
\end{equation}
is zero if and only if $T_{s\to t}$ is an exact gauge transformation
with trivial curvature on the language manifold.
\end{proposition}

\subsection{Translationese as Gauge Drift}

A well-documented phenomenon in MT is \textbf{translationese}~\cite{gellerstam1986translationese,
volansky2015translationese}: translated texts exhibit systematic
lexical, syntactic, and stylistic deviations from native texts in the
same language.

\begin{definition}[Translationese as Gauge Drift]
\label{def:translationese}
Let $x_t = T_{s\to t}(x_s)$ be a translation. The \textbf{translationese
degree} $\tau(x_t)$ is the divergence from the native gauge:
\begin{equation}
  \tau(x_t) = \KL(P_{\text{native}}(x) \| P_{\text{translated}}(x)),
\end{equation}
where $P_{\text{native}}$ is the distribution of native texts and
$P_{\text{translated}}$ is the distribution of translated texts in
the target language. Translationese is \textbf{gauge drift}---a
residual of the source language gauge that persists after the
gauge transformation $T_{s\to t}$.
\end{definition}

\begin{corollary}[Gauge Drift as Novely in Cercis]
\label{cor:translationese_cercis}
The \Cercis{} score between a translation and a native reference
captures translationese: a low $Q$ (agreement with native reference)
combined with detectable $N$ (systematic deviations) flags
translationese. This is exactly the \SCX{} novelty detection
mechanism applied to MT output.
\end{corollary}


% ================================================================
\section{Discussion}
\label{sec:discussion}
% ================================================================

\subsection{Summary of Contributions}

We have established a gauge-theoretic interpretation of machine
translation within the \SCX{} framework, providing nine exact
correspondences between MT linguistics and the \SCX{} audit formalism:

\begin{enumerate}[leftmargin=*, label=(\arabic*)]
  \item \textbf{Translation as gauge transformation} (Section~\ref{sec:translation_gauge}):
        source to target language is a coordinate chart transition that
        preserves gauge-invariant semantic content.
  \item \textbf{BLEU as Cercis} (Section~\ref{sec:bleu_cercis}):
        BLEU's n-gram precision matches the quality term $Q$, and its
        brevity penalty matches the novelty term $\eta N$.
  \item \textbf{Multilingual models as multi-expert systems}
        (Section~\ref{sec:multilingual}): shared latent space = multi-expert
        consensus space; cross-lingual transfer = gauge-invariant feature extraction.
  \item \textbf{Code-switching as gauge ambiguity} (Section~\ref{sec:codewitching}):
        alternating languages mid-utterance creates undefined gauge at switch points.
  \item \textbf{Universal Grammar as gauge-invariant substrate}
        (Section~\ref{sec:universal_grammar}): UG = principles (gauge constraints)
        + parameters (gauge degrees of freedom).
  \item \textbf{MT evaluation as audit} (Section~\ref{sec:evaluation_audit}):
        inter-annotator agreement = \Cercis{} score; MQM = multi-probe audit.
  \item \textbf{Word embeddings as gauge fields} (Section~\ref{sec:embeddings}):
        cross-lingual embedding alignment = gauge connection learning.
  \item \textbf{Neural MT as learned gauge fixing} (Section~\ref{sec:nmt_gauge}):
        transformer = gauge-equivariant architecture; attention = gauge transformation.
  \item \textbf{Translationese as gauge drift} (Section~\ref{sec:nmt_gauge}):
        systematic deviations from native texts are residual source-gauge effects.
\end{enumerate}

\subsection{Predictions and Testable Implications}

The gauge-theoretic framework yields quantitative predictions:

\begin{enumerate}[leftmargin=*, label=(\arabic*)]
  \item \textbf{BLEU--Cercis equivalence}. The BLEU score of a translation
        system can be interpreted as the log-\Cercis{} score, implying that
        the multiplicative structure of BLEU (product of precisions $\times$ BP)
        is the natural exponential-family structure of consensus scoring.
  \item \textbf{Gauge invariance of cross-lingual transfer}. Zero-shot
        transfer between languages $L_a$ and $L_b$ in multilingual models
        should be proportional to $1 / d_H(L_a, L_b)$, where $d_H$ is the
        Hausdorff distance in the shared representation space.
  \item \textbf{Code-switching perplexity bound}. The perplexity of a
        code-switched utterance should be strictly bounded below by the
        perplexity of its best monolingual paraphrase, providing a
        detection criterion for CS.
  \item \textbf{Translationese decay with model capacity}. As model capacity
        increases, the translationese degree $\tau(x_t)$ should decrease,
        approaching zero only in the infinite-capacity limit (perfect
        gauge fixing).
  \item \textbf{Embedding alignment and phylogenetic distance}. The alignment
        error $d_{\text{align}}(L_a, L_b)$ should correlate with the
        phylogenetic distance between $L_a$ and $L_b$. Language pairs
        within the same family should align better than cross-family pairs.
\end{enumerate}

\subsection{What This Paper Does Not Claim}

\begin{itemize}[leftmargin=*]
  \item \textbf{We do not claim} that language is literally a gauge field
        in the physicist's sense. The gauge analogy is a structural
        isomorphism at the level of formal properties---coordinate
        dependence, invariance, and parallel transport---not a claim
        about the physical substrate of language.
  \item \textbf{We do not claim} that multilingual models are explicitly
        computing gauge transformations. Rather, their learned representations
        \emph{realize} an approximate gauge structure as a byproduct of the
        MLM objective and the shared vocabulary.
  \item \textbf{We do not claim} that BLEU is identical to the \Cercis{}
        score. The mapping is at the level of functional form: both are
        precision-weighted geometric means with a novelty penalty. The
        numerical values differ because BLEU operates on n-gram overlap
        while \Cercis{} operates on expert agreement.
\end{itemize}

\subsection{Open Problems}

\begin{enumerate}[leftmargin=*, label=(\arabic*)]
  \item \textbf{Gauge-invariant MT metrics}. Can we construct MT metrics
        that are explicitly gauge-invariant---i.e., that measure translation
        quality without reference to a specific language coordinate chart?
        This would require a representation of meaning directly in the
        invariant space $\M$.
  \item \textbf{The gauge group of human languages}. What is the full
        structure of the gauge groupoid? Is it a Lie groupoid, and if so,
        what are its generators (the ``linguistic Lie algebra'')?
  \item \textbf{Unsupervised gauge fixing}. Can a model learn the gauge
        transformation $T_{s\to t}$ without parallel data, using only
        the gauge invariance constraint (e.g., cycle consistency)?
        This would be the unsupervised MT analog of unsupervised
        alignment~\cite{conneau2018word}.
  \item \textbf{Quantum effects in translation}. If translation gauge
        fields admit a Hamiltonian formulation, are there quantum-like
        effects---translation tunneling, superposition of translations,
        interference between translation paths?
  \item \textbf{The UG manifold dimension}. What is the intrinsic
        dimension $d_{\UG}$ of the Universal Grammar subspace in
        multilingual model representations? Does this correlate with
        the number of linguistic parameters ($\sim 20-50$ in
        generative grammar)?
  \item \textbf{Gauge anomaly cancellation}. Some linguistic properties
        appear to violate gauge invariance---e.g., the fact that
        grammatical gender does not translate across certain language
        pairs. Are these ``gauge anomalies'' that cancel when considering
        the full system?
  \item \textbf{Computational complexity of gauge fixing}. What is the
        computational complexity of learning the optimal gauge
        transformation $T_{s\to t}$? Does the gauge structure reduce
        the sample complexity of translation learning?
\end{enumerate}

\subsection{Conclusion}

We have presented a unified gauge-theoretic interpretation of machine
translation through the lens of the \SCX{} framework. The central
insight is that language---as a coordinate system for semantic
content---exhibits the same gauge structure as physical theories:
meaning is gauge-invariant, surface form is gauge-variant, and
translation is the gauge transformation between language coordinate
charts. This perspective connects BLEU scores, multilingual models,
code-switching, universal grammar, and MT evaluation under a single
formal framework, revealing deep structural analogies between
linguistics and the theory of multi-expert consensus.

The gauge-theoretic view is not merely descriptive---it generates
testable predictions about translation quality, cross-lingual
transfer, code-switching detection, and representation learning.
We hope that this perspective will inform the design of next-generation
MT systems that explicitly exploit gauge invariance, and will
deepen our understanding of the mathematical structure underlying
human language.


% ================================================================
% Appendix A: Glossary of Notation
% ================================================================
\section*{Appendix A: Glossary of Notation}
\label{app:notation}

\begin{table}[H]
\centering
\begin{tabular}{@{}ll@{}}
\toprule
\textbf{Symbol} & \textbf{Meaning} \\
\midrule
$\Lang_a = (\Vocab_a, \Phi_a, \Sigma_a)$ & Language coordinate system \\
$\Vocab_a$ & Vocabulary of language $L_a$ \\
$\Phi_a: \Vocab_a^* \to \M$ & Meaning mapping \\
$\Sigma_a \subset \Vocab_a^*$ & Grammatical sentences \\
$\M$ & Abstract meaning space \\
$\LangManifold$ & Language manifold (Riemannian) \\
$T_{a\to b}: \Sigma_a \to \Sigma_b$ & Translation gauge transformation \\
$\I(x)$ & Gauge-invariant content of sentence $x$ \\
$\Gamma$ & Linguistic connection (parallel transport) \\
$d_H(L_a, L_b)$ & Hausdorff distance between language clusters \\
\midrule
$p_n$ & N-gram precision of order $n$ \\
$\text{BP}(r, h)$ & Brevity penalty between reference and hypothesis \\
$\Cercis_{\text{ngram}}(r, h)$ & N-gram Cercis score \\
$\kappa$ & Inter-annotator agreement (Fleiss' kappa) \\
$S_{\text{MQM}}$ & MQM weighted Cercis score \\
\midrule
$A_a: \Vocab_a \to \R^d$ & Embedding gauge field \\
$W_{a\to b}$ & Cross-lingual alignment matrix \\
$d_{\text{align}}(L_a, L_b)$ & Alignment error distance \\
$\Al_{a\to b}$ & Gauge connection (alignment map) \\
\midrule
$\UG$ & Universal Grammar (gauge-invariant substrate) \\
$\theta_a$ & Parameter values of language $L_a$ \\
$G(x)$ & Gauge ambiguity score (CS detection) \\
$\tau(x_t)$ & Translationese degree (gauge drift) \\
\bottomrule
\end{tabular}
\end{table}


% ================================================================
% Appendix B: Numerical Verification
% ================================================================
\section*{Appendix B: Numerical Verification}
\label{app:numerical}

The Python script \texttt{verify\_linguistics.py} accompanying this
paper provides numerical verification of the key quantitative claims:

\begin{enumerate}[leftmargin=*, label=(\arabic*)]
  \item \textbf{BLEU as Cercis}: Verification that the BLEU score
        structurally decomposes into quality ($Q$ = n-gram precision)
        and novelty ($N$ = brevity penalty) terms.
  \item \textbf{Cross-lingual alignment}: Verification that embedding
        alignment via Procrustes learns a gauge connection, with
        P@k accuracy as a \Cercis{} score.
  \item \textbf{Code-switching detection}: Verification that CS
        utterances have lower \Cercis{} scores against monolingual
        references than monolingual utterances.
  \item \textbf{UG invariance}: Verification that multilingual model
        representations cluster by semantic content across language
        families, with cross-lingual Hausdorff distance as the
        invariance measure.
  \item \textbf{MQM audit simulation}: Verification that multi-annotator
        MT evaluation follows the \SCX{} audit protocol, with
        inter-annotator agreement as the \Cercis{} score.
  \item \textbf{Precision--recall decomposition}: Verification that
        different MT metrics correspond to different $Q$/$N$ weightings.
  \item \textbf{Transformer gauge equivariance}: Verification that
        transformer representations are approximately equivariant
        under embedding rotations.
\end{enumerate}

All numerical claims in the paper can be reproduced by running:
\begin{verbatim}
python verify_linguistics.py
\end{verbatim}
with \texttt{numpy} and \texttt{scipy} installed. Expected runtime
on a modern laptop: $\sim$ 30 seconds.


% ================================================================
% Appendix C: Proof of the BLEU--Cercis Isomorphism
% ================================================================
\section*{Appendix C: Proof of the BLEU--Cercis Isomorphism}
\label{app:bleu_cercis}

\begin{theorem}[Strong BLEU--Cercis Isomorphism]
\label{thm:strong_bleu_cercis}
\rigorFull
For any reference translation $r$ and hypothesis translation $h$
with $n$-gram statistics $p_n$ and brevity penalty BP, the log-BLEU
score is a \Cercis{} score $S(r, h)$ with quality and novelty terms:
\begin{equation}
  \log\BLEU(r, h) = Q(r, h) + \eta(h) \cdot N(r, h),
\end{equation}
where the mapping is bijective up to a multiplicative constant.
Conversely, any \Cercis{} score $S$ with $K$ experts and binary
quality/novelty decomposition can be expressed as a BLEU score
with $N = K$ n-gram orders and a suitably defined brevity penalty.
\end{theorem}

\begin{proof}
\noindent\textbf{Step 1: Forward direction (BLEU $\to$ Cercis).}
By Theorem~\ref{thm:bleu_cercis}, we have the identification:
\begin{equation}
  Q(r, h) = \frac{1}{N}\sum_{n=1}^N \log p_n, \quad
  \eta(h) \cdot N(r, h) = \log\text{BP}(r, h).
\end{equation}
The mapping is injective: different BLEU scores map to different
pairs $(Q, N)$ because the precision values $p_n$ are uniquely
determined by the n-gram counts.

\smallskip
\noindent\textbf{Step 2: Reverse direction (Cercis $\to$ BLEU).}
Given a \Cercis{} score $S(s) = Q(s) + \eta(t)N(s)$ over $K$
experts, define:
\begin{align}
  p_n &= \exp(Q_n) \quad \text{(where } Q_n \text{ is the per-expert quality)},\\
  \text{BP} &= \exp(\eta(t)N(s)).
\end{align}
Set $N = K$ and construct $p_n$ from $Q_n$. This defines a BLEU
score: $\BLEU = \text{BP} \cdot (\prod_{n=1}^K p_n)^{1/K} = e^{S(s)}$.
The mapping is surjective onto the range of \Cercis{} scores with
the binary Q/N decomposition.

\smallskip
\noindent\textbf{Step 3: Compositionality.}
The isomorphism respects composition:
\begin{equation}
  \BLEU(r, h_1 \oplus h_2) = \BLEU(r, h_1) \cdot \BLEU(r, h_2)
  \iff \Cercis(r, h_1 \oplus h_2) = \Cercis(r, h_1) + \Cercis(r, h_2),
\end{equation}
where $\oplus$ denotes concatenation of hypotheses.
\end{proof}


% ================================================================
\begin{thebibliography}{99}

\bibitem{papineni2002bleu}
K.~Papineni, S.~Roukos, T.~Ward, and W.-J.~Zhu,
\emph{BLEU: A method for automatic evaluation of machine translation},
Proc.\ ACL (2002), 311--318.

\bibitem{devlin2019bert}
J.~Devlin, M.-W.~Chang, K.~Lee, and K.~Toutanova,
\emph{BERT: Pre-training of deep bidirectional transformers for
  language understanding},
Proc.\ NAACL-HLT (2019), 4171--4186.

\bibitem{conneau2020unsupervised}
A.~Conneau, K.~Khandelwal, N.~Goyal, V.~Chaudhary, G.~Wenzek,
  F.~Guzman, E.~Grave, M.~Ott, L.~Zettlemoyer, and V.~Stoyanov,
\emph{Unsupervised cross-lingual representation learning at scale},
Proc.\ ACL (2020), 8440--8451.

\bibitem{hassan2018achieving}
H.~Hassan, A.~Aue, C.~Chen, V.~Chowdhary, J.~Clark, C.~Federmann,
  X.~Huang, M.~Junczys-Dowmunt, W.~Lewis, M.~Li, S.~Liu, T.-Y.~Liu,
  R.~Luo, A.~Menezes, T.~Qin, F.~Seide, X.~Tan, F.~Tian, L.~Wu,
  S.~Wu, Y.~Xia, D.~Zhang, Z.~Zhang, and M.~Zhou,
\emph{Achieving human parity on automatic Chinese to English translation},
arXiv:1803.05567 (2018).

\bibitem{vaswani2017attention}
A.~Vaswani, N.~Shazeer, N.~Parmar, J.~Uszkoreit, L.~Jones,
  A.~N.~Gomez, L.~Kaiser, and I.~Polosukhin,
\emph{Attention is all you need},
Advances in Neural Information Processing Systems 30 (2017), 5998--6008.

\bibitem{mikolov2013distributed}
T.~Mikolov, I.~Sutskever, K.~Chen, G.~S.~Corrado, and J.~Dean,
\emph{Distributed representations of words and phrases and their
  compositionality},
Advances in Neural Information Processing Systems 26 (2013), 3111--3119.

\bibitem{pennington2014glove}
J.~Pennington, R.~Socher, and C.~D.~Manning,
\emph{GloVe: Global vectors for word representation},
Proc.\ EMNLP (2014), 1532--1543.

\bibitem{bojanowski2017enriching}
P.~Bojanowski, E.~Grave, A.~Joulin, and T.~Mikolov,
\emph{Enriching word vectors with subword information},
Trans.\ ACL 5 (2017), 135--146.

\bibitem{sennrich2016improving}
R.~Sennrich, B.~Haddow, and A.~Birch,
\emph{Improving neural machine translation models with monolingual data},
Proc.\ ACL (2016), 86--96.

\bibitem{conneau2018word}
A.~Conneau, G.~Lample, M.~Ranzato, L.~Denoyer, and H.~Jegou,
\emph{Word translation without parallel data},
Proc.\ ICLR (2018).

\bibitem{artetxe2018massively}
M.~Artetxe, G.~Labaka, and E.~Agirre,
\emph{A robust self-learning method for fully unsupervised
  cross-lingual mappings of word embeddings},
Proc.\ ACL (2018), 789--798.

\bibitem{pires2019multilingual}
T.~Pires, E.~Schlinger, and D.~Garraete,
\emph{How multilingual is multilingual BERT?},
Proc.\ ACL (2019), 4996--5001.

\bibitem{wu2024bilingual}
S.~Wu and M.~Dredze,
\emph{Bilingual alignment and cross-lingual transfer},
Trans.\ ACL 12 (2024), 101--118.

\bibitem{kudo2018sentencepiece}
T.~Kudo and J.~Richardson,
\emph{SentencePiece: A simple and language independent subword
  tokenizer and detokenizer for neural text processing},
Proc.\ EMNLP (2018), 66--71.

\bibitem{lin2004rouge}
C.-Y.~Lin,
\emph{ROUGE: A package for automatic evaluation of summaries},
Proc.\ ACL Workshop on Text Summarization (2004).

\bibitem{banerjee2005meteor}
S.~Banerjee and A.~Lavie,
\emph{METEOR: An automatic metric for MT evaluation with improved
  correlation with human judgments},
Proc.\ ACL Workshop on Intrinsic and Extrinsic Evaluation Measures
  for MT (2005), 65--72.

\bibitem{rei2020comet}
R.~Rei, C.~Stewart, A.~C.~Farinha, and A.~Lavie,
\emph{COMET: A neural framework for MT evaluation},
Proc.\ EMNLP (2020), 2685--2702.

\bibitem{myers2002code}
C.~Myers-Scotton,
\emph{Contact Linguistics: Bilingual Encounters and Grammatical Outcomes},
Oxford University Press, 2002.

\bibitem{chomsky1965aspects}
N.~Chomsky,
\emph{Aspects of the Theory of Syntax},
MIT Press, 1965.

\bibitem{chomsky1981lectures}
N.~Chomsky,
\emph{Lectures on Government and Binding},
Foris, 1981.

\bibitem{greenberg1963universals}
J.~H.~Greenberg,
\emph{Universals of Language},
MIT Press, 1963.

\bibitem{dryer2013wals}
M.~S.~Dryer and M.~Haspelmath (eds.),
\emph{The World Atlas of Language Structures Online},
Max Planck Institute for Evolutionary Anthropology, 2013.

\bibitem{lommel2014multidimensional}
A.~Lommel, A.~G{\'o}rska, and H.~Meliksetian,
\emph{Multidimensional quality metrics (MQM): A framework for
  declaring and describing translation quality metrics},
Tradulex (2014), 455--463.

\bibitem{castilho2017guide}
S.~Castilho, J.~Moorkens, F.~Gaspari, I.~Calleja, C.~Way,
  and A.~Panea,
\emph{A guide to evaluating machine translation quality},
The Routledge Handbook of Translation and Technology (2017), 137--156.

\bibitem{gellerstam1986translationese}
M.~Gellerstam,
\emph{Translationese in Swedish novels translated from English},
Translation Studies in Scandinavia (1986), 88--95.

\bibitem{volansky2015translationese}
V.~Volansky, N.~Ordan, and S.~Wintner,
\emph{On the features of translationese},
Digital Scholarship in the Humanities 30 (2015), 98--118.

\bibitem{ba2016layer}
J.~L.~Ba, J.~R.~Kiros, and G.~E.~Hinton,
\emph{Layer normalization},
arXiv:1607.06450 (2016).

\bibitem{scx_theory}
SCX,
\emph{State-Conditioned Expertise: A Complete Theory of Label Noise
  Detection via Multi-Expert Consistency},
SCX Technical Report, 2026.

\bibitem{scx_spring}
SCX,
\emph{Spring: A Self-Evolving Gatekeeper with Provable Convergence},
SCX Technical Report, 2026.

\bibitem{scx_compactness}
SCX,
\emph{Compactness and SCX: From Finite to Infinite Audit Transfer},
SCX Technical Report, 2026.

\bibitem{scx_qg_audit}
SCX,
\emph{Quantum Gravity Audit Equivalence: Gauge-Equivalent Descriptions
  and the Compactness-Inseparability Criterion},
SCX Technical Report, 2026.

\end{thebibliography}

\end{document}
